% !TeX root = ../main.tex

\sysusetup{
  keywords  = {Agent Skill；价值发现；SkillRank；条件排序；语义建模；云端系统},
  keywords* = {Agent Skill; value discovery; SkillRank; conditional ranking; semantic modeling; cloud system},
}

\begin{abstract}
  随着大语言模型驱动的智能体（Agent）技术快速普及，围绕 Agent Skill 的开发与分发生态正在高速扩张。Skill 作为 Agent 能力扩展的核心载体，本质上是一个以 Description 为键、以 Content 为值的上下文注入单元。然而，随着可用 Skill 数量从数十增长至数千乃至数万，Agent 的意图匹配准确率急剧下降——实验表明，当本地 Skill 数量超过 200 时，准确率可跌至约 20\%。与此同时，现有 Skill 平台缺乏系统化的质量评估机制，高价值 Skill 淹没于大量低质重复内容中，普通用户无从判断应启用哪些 Skill。

  针对上述问题，本文设计并实现了面向大规模 Agent Skill 价值发现的云端系统 Taste，其核心算法为 SkillRank。本文将 Skill 形式化为二元组 $\langle Description, Content \rangle$，其中 Description 决定 Skill 在需求空间中的语义坐标，Content 决定 Skill 的实质质量。基于这一建模，本文提出 Skill 价值评估是一个\textbf{条件排序问题}：仅在语义邻域（相同描述聚类）内对 Content 质量排序才有意义。SkillRank 借鉴 PageRank 的图传播思想，以 Content 语义相似度构建稀疏 K-NN 亲缘图，在其上运行 PageRank 迭代，使被广泛模仿的解法模式获得更高分值，并通过 FAISS 近似最近邻搜索将复杂度从 $O(n^2)$ 降至 $O(n \log n)$。

  在评估方面，本文以从 ClawHub 平台爬取的 2,103 条含 GitHub Star 数据的真实 Skill 为数据集，以聚类内 GitHub Star 排序为 Ground Truth，采用 Spearman $\rho$ 作为评测指标。经 AutoResearch 风格的自动化迭代优化，最优配置（$K=15$，damping factor$=0.82$）下 Spearman $\rho$ 达到 0.71，相比随机基线的 0.23 提升显著；在规模为 10,000 条 Skill 时，相比朴素 $O(n^2)$ 方案加速约 40 倍。

  系统层面，Taste 提供完整的云端 Skill 发现与使用平台，后端采用 Python FastAPI 与 PostgreSQL，前端提供基于 Plotly.js 的三维语义地形图可视化（XY 为 Description embedding 降维坐标，Z 为 SkillRank 分），支持用户输入 context 与参数 K 后直接获取最优 Skill 推荐集合。
\end{abstract}

\begin{abstract*}
  As large language model-driven agent technology rapidly proliferates, the ecosystem surrounding Agent Skills is expanding at scale. A Skill serves as the core unit of agent capability extension, fundamentally acting as a context injection dictionary entry with Description as key and Content as value. However, as the number of available Skills grows from dozens to thousands or even tens of thousands, the agent's intent-matching accuracy drops sharply — experiments show that accuracy falls to approximately 20\% when the local Skill count exceeds 200. Meanwhile, existing Skill platforms lack systematic quality evaluation mechanisms, and high-value Skills are buried among large amounts of low-quality, redundant content, leaving ordinary users unable to determine which Skills to enable.

  To address these problems, this thesis designs and implements Taste, a cloud-based system for large-scale Agent Skill value discovery, with SkillRank as its core algorithm. Skills are formalized as a dual tuple $\langle Description, Content \rangle$, where Description determines a Skill's semantic coordinates in the requirement space, and Content determines its substantive quality. Based on this modeling, Skill value evaluation is framed as a \textbf{conditional ranking problem}: ranking Content quality is meaningful only within a semantic neighborhood (the same description cluster). SkillRank adopts the graph propagation idea of PageRank, constructs a sparse K-NN affinity graph using Content semantic similarity, runs PageRank iteration on it so that widely-imitated solution patterns receive higher scores, and reduces complexity from $O(n^2)$ to $O(n \log n)$ via FAISS approximate nearest neighbor search.

  For evaluation, a real-world dataset of 2,103 Skills with GitHub Star counts crawled from the ClawHub platform is used, with within-cluster GitHub Star ranking as the Ground Truth and Spearman $\rho$ as the metric. Through AutoResearch-style automated iterative optimization, the best configuration ($K=15$, damping factor $=0.82$) achieves a Spearman $\rho$ of 0.71, a significant improvement over the random baseline of 0.23; at a scale of 10,000 Skills, it achieves approximately 40x speedup over the naive $O(n^2)$ approach.

  At the system level, Taste provides a complete cloud-based Skill discovery and consumption platform. The backend uses Python FastAPI and PostgreSQL, and the frontend provides a 3D semantic terrain visualization based on Plotly.js (XY as reduced-dimension Description embedding coordinates, Z as SkillRank score), allowing users to input a context and parameter K to directly obtain an optimal Skill recommendation set.
\end{abstract*}
